\documentclass[aspectratio=169, 11pt]{beamer}
\usetheme{Madrid}
\usecolortheme{whale}

\usepackage[T1]{fontenc}
\usepackage{lmodern}
\usepackage{graphicx}
\usepackage{booktabs}
\usepackage{hyperref}
\usepackage{xcolor}
\usepackage{amsmath}

% --- Placeholders (filled by the skill) ---
\newcommand{\casedate}{YYYY-MM-DD}
\newcommand{\casetopic}{Topic Name}
\newcommand{\caseslug}{topic_slug}

\title{Today's Market Case}
\subtitle{\casetopic{} --- \casedate}
\author{EF5342 --- Financial Systems, Markets and Instruments}
\institute{City University of Hong Kong}
\date{\casedate}

\begin{document}

%=== Slide 1: Title ===%
\begin{frame}
\titlepage
\end{frame}

%=== Slide 2: Context + Key Numbers ===%
\begin{frame}{What Happened: [EVENT HEADLINE]}
[2--3 sentence summary of the market event. What happened, why it matters,
and what markets are reacting to.]

\vspace{1em}
\begin{table}
\centering
\small
\begin{tabular}{lrr}
\toprule
Metric & Value & Change \\
\midrule
Indicator 1 & Level & \% or bps change \\
Indicator 2 & Level & \% or bps change \\
Indicator 3 & Level & \% or bps change \\
Indicator 4 & Level & \% or bps change \\
\bottomrule
\end{tabular}
\end{table}

\vspace{0.5em}
{\footnotesize \textit{Source:} [News source], [Date]. Data: Yahoo Finance.}
\end{frame}

%=== Slide 3: Chart ===%
\begin{frame}{[CHART TITLE, e.g., ``10-Year Treasury Yield, 3-Month Trailing'']}
\begin{center}
% The skill replaces the \fbox below with \includegraphics[width=0.85\textwidth]{chart.png}
\fbox{\parbox{0.85\textwidth}{\centering\vspace{2cm}[CHART PLACEHOLDER]\vspace{2cm}}}
\end{center}
\vspace{-0.5em}
{\footnotesize [One-line caption describing what the chart shows.]}
\end{frame}

%=== Slide 4: Historical Parallel + Theory Link ===%
\begin{frame}{Putting It in Context}
\textbf{Historical parallel:}
[When did something similar happen? 1--2 sentences comparing today's event to
a past episode, noting what is similar and what is different.]

\vspace{1em}
\textbf{Connection to [WEEK TOPIC]:}
[1--2 sentences linking the event to the course concept being taught that week.
If there is a key formula or diagram, include it here, e.g.:]

\[
\text{Yield to Maturity: } P = \sum_{t=1}^{T} \frac{C}{(1+y)^t} + \frac{F}{(1+y)^T}
\]
\end{frame}

%=== Slide 5: Discussion Questions ===%
\begin{frame}{Discussion Questions}
% [Bloom: Apply]
\textbf{Q1.} [First question --- directly applies the week's concept to the event.
Ask students to use a framework, formula, or model to explain what happened.]

\vspace{1em}
% [Bloom: Analyse]
\textbf{Q2.} [Second question --- compare or contrast with the historical parallel,
or analyse cause-and-effect using course concepts.]

\vspace{1em}
% [Bloom: Evaluate]
\textbf{Q3.} [Third question --- judgment-oriented. Ask students to take a position
on a policy or investment decision, justifying with course concepts.]
\end{frame}

\end{document}
